@ID: OW26D1P1I
@BIDANG: Kombinatorika
Tentukan banyaknya cara menempatkan $n$ bola berbeda ke dalam $n$ kotak berbeda sehingga tepat satu kotak kosong.

@ID: OW26D1P2I
@BIDANG: Analisis Real
Misalkan
$$
\sum_{n=1}^{\infty} a_n
$$
adalah deret dengan suku-suku positif. Jika

$$
\lim_{n\to\infty}\frac{a_{n+1}}{a_n}
=
\frac{a^2+5}{6},
$$

maka tentukan himpunan semua nilai $a$ yang menjamin deret tersebut konvergen.

@ID: OW26D1P3I
@BIDANG: Analisis Real
Diberikan barisan fungsi $f_n:[0,1]\to\mathbb{R}$ dengan

$$
f_n(x):=n^2x^2(1-x)^n.
$$

Tentukan nilai

$$
\lim_{n\to\infty}
\sup_{0\le x\le1}
f_n(x).
$$

@ID: OW26D1P4I
@BIDANG: Struktur Aljabar
Jika

$$
A:=\{\,n\in\mathbb{Z}_{1400}\mid
n\not\equiv3\pmod8,\;
n\not\equiv4\pmod{25},\;
n\not\equiv5\pmod7
\,\},
$$

tentukan nilai $|A|$.

@ID: OW26D1P5I
@BIDANG: Struktur Aljabar
Jika

$$
\alpha:=(1\,2\,3\,4\,5\,6)(1\,2\,3)(1\,4\,5)(1\,6\,5\,4\,3\,2)\in S_7,
$$

tentukan nilai $\alpha^{99}.$

@ID: OW26D1P1U
@BIDANG: Analisis Real
Tentukan bilangan real terkecil $c$ sehingga untuk setiap $x>0$ berlaku

$$
\ln(1+2e^{3x})<4x+c.
$$

@ID: OW26D1P2U
@BIDANG: Analisis Real
Misalkan fungsi $f:[a,b]\to\mathbb{R}$ monoton naik dan $g:[a,b]\to\mathbb{R}$ terintegral pada $[a,b]$. Jika $g(x)\ge0$ untuk setiap $x\in[a,b]$, buktikan terdapat $c\in[a,b]$ sehingga

$$
\int_a^b f(x)g(x)\,dx
=
f(a)\int_a^c g(x)\,dx
+
f(b)\int_c^b g(x)\,dx.
$$

@ID: OW26D1P3U
@BIDANG: Struktur Aljabar
Diberikan suatu ring $(R,+,\times)$ dengan identitas perkalian $1_R$. Jika untuk setiap $r\in R$, minimal salah satu dari $r$ atau $1_R+r$ memiliki invers terhadap operasi perkalian, buktikan bahwa ring $R$ memiliki tepat satu ideal maksimal.

@ID: OW26D1P4U
@BIDANG: Struktur Aljabar
Misalkan $G$ adalah himpunan tak kosong dengan operasi biner asosiatif, ditulis sebagai perkalian. Misalkan untuk setiap $x,y,z\in G$ berlaku

$$
xy=zxyz.
$$

Buktikan:

(a) Operasi pada $G$ komutatif.

(b) Jika

$$
H:=\{xy\mid x,y\in G\},
$$

maka $H$ adalah grup terhadap operasi yang sama.

@ID: OW26D1P5U
@BIDANG: Kombinatorika
Buktikan bahwa $(4n)!$ habis dibagi oleh $2^{3n}3^n$ untuk setiap bilangan asli $n$.













@ID: OW26D2P1I
@BIDANG: Aljabar Linear
Misalkan

$$
X=
\begin{pmatrix}
2&0\\
2&6
\end{pmatrix}.
$$

Tentukan matriks $Y$ berukuran $2\times2$ yang memenuhi persamaan

$$
2XY=21X-2Y.
$$

@ID: OW26D2P2I
@BIDANG: Analisis Kompleks
Apabila diketahui semua akar persamaan

$$
\left(\frac{1+iz}{1-iz}\right)^{2026}
=
\frac12+ik
$$

merupakan bilangan real, tentukan nilai terbesar $k\in\mathbb{R}$ yang memenuhi.

@ID: OW26D2P3I
@BIDANG: Analisis Kompleks
Untuk $|z-2i|<2$, didefinisikan

$$
f(z):=
\sum_{n=1}^{\infty}
\frac{(-1)^{n-1}}{n}
\left(\frac{z-2i}{2i}\right)^n.
$$

Tentukan nilai

$$
|f'(1+i)|^2.
$$

@ID: OW26D2P4I
@BIDANG: Aljabar Linear
Misalkan $A=(a_{ij})$ dan $B=(b_{ij})$ matriks berukuran $2026\times2026$ dengan

$$
a_{ij}:=
\begin{cases}
i+j,&\text{jika } i\le j,\\
0,&\text{jika } i>j,
\end{cases}
$$

dan

$$
b_{ij}:=
\begin{cases}
\dfrac{1}{i+j},&\text{jika } i\ge j,\\
0,&\text{jika } i<j.
\end{cases}
$$

Tentukan nilai

$$
\operatorname{tr}(AB).
$$

Catatan: $\operatorname{tr}(M)$ menyatakan trace dari matriks $M$.

@ID: OW26D2P5I
@BIDANG: Kombinatorika
Tentukan banyaknya bilangan 10-digit yang seluruh digitnya tidak boleh nol dan setiap bilangan yang dibentuk oleh dua angka berdekatan dapat dibagi oleh salah satu dari $13$, $17$, $23$, atau $37$.

@ID: OW26D2P1U
@BIDANG: Analisis Kompleks
Tentukan semua bilangan kompleks yang memenuhi persamaan

$$
z^{2026}-z^{2025}+z^{2024}-\cdots-z+1=0.
$$

@ID: OW26D2P2U
@BIDANG: Analisis Kompleks
Misalkan fungsi kompleks $f$ analitik pada domain $C$. Buktikan bahwa fungsi kompleks

$$
g(z):=\overline{f(\overline{2026z})}
$$

juga analitik pada $C$.

@ID: OW26D2P3U
@BIDANG: Aljabar Linear
Misalkan $P$ adalah ruang semua polinom atas $\mathbb{R}$. Didefinisikan transformasi linear $T:P\to P$ dengan

$$
T(f(x)):=xf(x)
$$

untuk setiap polinom $f(x)$ di $P$.

(a) Tunjukkan bahwa $T$ tidak memiliki nilai eigen real.

(b) Apakah terdapat suatu transformasi linear $S:P\to P$ sehingga $T\circ S$ merupakan pemetaan identitas? Jelaskan jawaban Anda.

@ID: OW26D2P4U
@BIDANG: Aljabar Linear
Misalkan $A$ suatu matriks berukuran $4\times4$ dengan $\operatorname{rank}(A)=2$. Buktikan bahwa terdapat vektor-vektor kolom $u_1,u_2,w_1,w_2$ di $\mathbb{R}^4$ sehingga

$$
u_1^Tu_2=0
$$

dan

$$
A=u_1w_1^T+u_2w_2^T.
$$

@ID: OW26D2P5U
@BIDANG: Kombinatorika
Buktikan bahwa

$$
\sum_{k=0}^{n}
(-1)^k
\binom{n}{k}
f(k)=0
$$

untuk setiap polinom $f$ dengan derajat kurang dari $n$.